% !TEX root = ../main.tex
\section{Impact/Financial Benefits}
\label{sec:impact}

\plantodo{S9, S12 (S16 bullet added this session)}{Quantify the impact of
the reference architecture and its optimization-demonstration capability.
Translate the S9 conflict analysis and S12 demonstration results into
concrete operational terms wherever possible (e.g., delay-minutes,
fuel burn, CO2, gate/runway utilization) and into US-dollar terms where a
credible conversion exists (e.g., published cost-per-delay-minute figures
from the literature review, Section~\ref{sec:litreview}). Also discuss the
ethical context and any unintended consequences of the recommendations
(equity-minded engineering: whose costs get externalized, and onto whom --
this connects directly to S8's "objectives whose costs are externalized
onto another stakeholder" and S9's conflict catalog).}

\textbf{Scope flag:} the 503 template expects a quantified financial-benefit
statement (the template's sample statements are dollar-savings claims from
industry DMAIC-style projects). This capstone is an architecture/SoS
project, not a cost-reduction project with a "before" baseline in dollars,
so a literal dollar-savings figure may not be available or meaningful.
Project\_To-Do List.md did not previously have a task that produces a
quantified/monetized impact figure -- an explicit "Impact / quantified
benefits" bullet was added to S16 this session so at least one deliverable
targets this section directly. Recommend framing impact primarily in
operational/systems terms (delay, fuel, emissions, predictability,
tradeoff-surface shape) with a dollar estimate offered as an illustrative,
clearly-caveated secondary figure rather than the headline metric.

% ---- Placeholders added 2026-10-09 (AI-drafted scaffolding, not report prose).
\subsection{Emissions: What They Mean to the Airline and to Society}
\label{sec:emissions-impact}

\plantodo{S9, S12}{\textbf{Airline level: put CO2 into dollars.} Take the
fuel difference between cases from Table~\ref{tab:outcomes}, convert to kg
of CO2, then to dollars. State the route used and its source. Candidates:
(1) the price of a carbon credit or allowance, if a scheme applies to the
flight; (2) passenger value, meaning revenue gained or lost because some
passengers choose the flight listed with lower emissions. Make the point
that fuel is already counted once as fuel cost and that pricing its CO2
counts the same kilograms again, on purpose. SOURCES NOT YET HELD.}

\plantodo{S9, S12}{\textbf{Passenger value.} Discuss how passengers see a
flight's emissions when booking, and whether that reaches the airline as
demand. The author's own practice of choosing lower-emission listings is the
motivation, not the evidence: cite a study. Note the gap between one
decision in one sector and the label a passenger saw at booking; the link is
the airline's average over many flights.}

\plantodo{S9, S12}{\textbf{System-of-systems level: keep CO2 in kg.} Report
total CO2 across all four aircraft for each case. Discuss what it means for
society and for how passengers perceive air travel. A social cost of carbon
may be quoted for scale, labelled as a cost nobody in the scenario pays.
Connect to S8's objectives whose costs are externalized onto another
stakeholder.}

\plantodo{S12}{\textbf{Scale it honestly.} One conflict in one sector is a
small number. If a per-year or fleet-wide figure is offered, state the
multiplier and its source (for example, how often level crossing conflicts
occur), and label it illustrative.}

\plantodo{S9}{\textbf{Who pays.} For the option the system prefers, name the
aircraft or airline that is worse off and by how much. This is the
equity discussion the template asks for.}
% ---- end placeholders 2026-10-09
